\documentclass[10pt,a4paper]{amsart}
\usepackage[margin=19mm]{geometry}
\usepackage{amsmath,amssymb,mathtools}
\usepackage[scr=rsfs,cal=euler]{mathalfa}
\usepackage{xfrac}
\usepackage[unicode,hidelinks,bookmarks=false]{hyperref}
\newcommand{\ie}{i.e.\ }
\newcommand{\cf}{cf.\ }
\newcommand{\resp}{respectively\ }
\newcommand{\Subsection}{\subsection}
\providecommand{\nobreakdash}{\nobreak\hbox{-}}
\setlength{\emergencystretch}{2em}
\multlinegap0pt
\usepackage{bm}
\usepackage{bbm}
\usepackage{dsfont}
\usepackage{xspace}
\usepackage{cancel}
\usepackage{mathtools}
\usepackage{amsthm}
\makeatletter
\@ifundefined{thm}{\newtheorem{thm}{Thm}}{}
\@ifundefined{prop}{\newtheorem{prop}[thm]{Proposition}}{}
\makeatother
\newcommand{\Trd}{\operatorname{Trd}}
\newcommand{\iso}{\cong}
\title[Extending degree-2 Azumaya algebras with C2-actions and exam]{Extending degree-2 Azumaya algebras with C2-actions and examples from character varieties of knot group}
\author{Justin Lawrence, Nicholas Rouse, Ben Williams}
\date{Source: arXiv:2510.25149v1 (CC BY 4.0)}
\begin{document}
\maketitle

\noindent\emph{Source and licence.} Extending degree-2 Azumaya algebras with C2-actions and examples from character varieties of knot group. Version arXiv:2510.25149v1 (CC BY 4.0), distributed under the Creative Commons Attribution 4.0 International licence, as recorded in the arXiv licence metadata for this exact version and shown on the abstract page https://arxiv.org/abs/2510.25149v1. Full rights notice: licenses/arxiv-2510.25149-CC-BY-4.0.txt.\par\smallskip

\noindent\textit{Editorial scope.} Counted unit: the Prop whose source label is \texttt{pr:symbolAlgebra} in the article (source0.tex lines 484--493, source label \texttt{pr:symbolAlgebra}), delivered together with the article's own complete original proof of it (source0.tex lines 495--535, one \texttt{proof} environment of 41 lines). No mathematics is written, repaired or re-derived: statement and proof are the article's own text line for line. Required context retained uncounted: none. Every definition, display and label of those lines is retained unchanged. Omitted from this package and not redistributed: 1--483, 494--494, 536--862. Nothing inside a retained excerpt is omitted or altered: every retained body is delivered exactly as the article's lines are written, so no omission receipt of any kind accompanies this package. Citations: the retained excerpts carry no citation command at all, so no bibliography entry of the article is transcribed and no reference is redistributed. Reference adaptation: none was needed and none was made; every reference inside the delivered unit resolves inside it, so no unresolved reference or citation is produced. Printed slips kept literal: no printed word, symbol, hypothesis or number of the counted statement or of its proof was altered. Layout and class: the article's own class is \texttt{amsart} with packages including geometry, hyperref, biblatex, breakurl, fourier, graphicx, tikz, amssymb, amsfonts, amsmath, mathtools, enumitem, xcolor, ifdraft; the wrapper is the lane's pinned \texttt{amsart} editorial wrapper at 10pt on A4, which restates the theorem-like environments the retained text needs and adds the article's own macro definitions that the retained text needs (\texttt{\textbackslash Trd}, \texttt{\textbackslash iso}), with extra wrapper packages (bm, bbm, dsfont, xspace, cancel, mathtools). The delivered numbering is the unit's own. Assets: no retained span contains a figure environment, a graphics-inclusion command, a PDF-page inclusion, a table or a TikZ picture; no image file is copied into this package, the package declares no asset dependency and no image file is redistributed with it. Licence: version arXiv:2510.25149v1 of this article is distributed under the Creative Commons Attribution 4.0 International licence, as recorded in the arXiv licence metadata for this exact version and shown on the abstract page https://arxiv.org/abs/2510.25149v1. A full rights notice is delivered at licenses/arxiv-2510.25149-CC-BY-4.0.txt. Characters: every retained excerpt is pure ASCII, so the delivered engine, XeLaTeX with its default Latin Modern text font, reports no missing character.
\begin{prop} \label{pr:symbolAlgebra}
 Suppose that $m_g$ and $m_h$ generate $A$ as an algebra, and that $I_g^2
  \neq 4$. Define
  \[ i = 2m_g  - I_g, \quad j' = 2m_h - I_h, \quad j= -i\left(j' - \frac{\langle i,j' \rangle}{\langle i , i \rangle }
      i \right).\]
  Then $i,j$ generate $A$ as an $F$-algebra, $ij=-ji$ and
  \[ i^2 = I_g^2 - 4, \quad j^2 = -(I_g^2 -4)(I_h^2-4) + (2I_{gh} - I_gI_h)^2. \]
  In particular, $A$ is a symbol algebra given by
  \[ A \iso \left( \frac{I_g^2 - 4, \, -(I_g^2 -4)(I_h^2-4) + (2I_{gh} - I_gI_h)^2}{F} \right).\]  
\end{prop}

\begin{proof}
Starting with the generators $m_g, m_h$ of $A$, we may take the traceless elements
\[ i = 2m_g - I_g,\qquad j'= 2 m_h -  I_h, \]
and then orthogonalize with respect to the bilinear form:
\[ j'' = j' - \frac{ \langle i, j' \rangle}{\langle i, i \rangle} i. \]
% With this definition, Proposition \ref{pr:anticommute} says that the traceless elements $i$ and $j''$, and they generate the $F$-algebra $A$. Note that
% \[ i^2 = 4 m_g^2 - 4I_g m_g + I_g^2 = I_g^2 - 4 \neq 0 \]
% so that $a$ is a unit.

\noindent We replace $j''$ by
\[ j = -i j''. \]
Together, $i, j$ are generators of the $F$-algebra $A$. Since $i$ and $j''$ are orthogonal and $j''$ is traceless,  $i$ and $-j'' = \bar \jmath''$ are also orthogonal. We calculate
\[ \Trd(j) = -ij'' + (-\bar \jmath'' \bar \imath) = -ij'' -j''i =  -2 \langle i, j'' \rangle = 0. \]
Furthermore,
\[ ij = -i^2 j'' = ij''i = -ji. \]
so that $i,j$ anticommute, i.e., are themselves orthogonal.

Therefore, $A \iso \left( \frac{ i^2, j^2 }{F} \right)$. It remains to calculate $i^2$ and $j^2$. One of these is easily written down directly:
 \[ i^2 = 4 m_g^2 - 4I_g m_g + I_g^2 = I_g^2 - 4. \]

Calculating $j^2$ is considerably more time-consuming. To simplify the notation, let us write
\[ Q = \frac{ \langle i, j' \rangle}{\langle i, i \rangle}. \] We may calculate explicitly

\begin{equation}
  \begin{aligned}  Q & = \frac{\Trd(i \bar \jmath')}{\Trd(i\bar \imath)} = \frac{\Trd(ij')}{\Trd(i^2)}\\
    &= \frac{\Trd(4m_gm_h -2m_gI_h -2m_hI_g +I_gI_h)}{\Trd(4m_g^2 -4m_gI_g +I^2_g) } \\
    &= \frac{2I_{gh} - I_g I_h}{I_g^2-4}.
  \end{aligned}
  \label{eq:17}
\end{equation}

\noindent Using this, we may write
\begin{align*} j^2 &= -(I_g^2-4)(b'')^2= (I_g^2-4)[(I_h^2 -4) - Q(ij' + j'i) + Q^2 (I_g^2-4)] \\
                   &= -(I_g^2-4)\left[(I_h^2 -4) -
                     \frac{2I_{gh}-I_gI_h}{I_g^2-4}(\Trd(4m_{gh}-2I_gm_h-2I_hm_g+I_gI_h))+ \right. \\
                     & \qquad \left.+\frac{(2I_{gh}-I_gI_h)^2}{(I_g^2-4)^{2}}(I_g^2-4)\right] \\
                   &= -(I_g^2-4)(I_h^2-4) +2(2I_{gh}-I_gI_h)^2-(2I_{gh}-I_gI_h)^2 \\
                   &=-(I_g^2-4)(I_h^2-4)  + (2I_{gh}-I_gI_h)^2.
\end{align*}
This completes the proof.
\end{proof}

\end{document}
